\section{Introduction}
%%%%% Field gap %%%%% 

Representations are \textit{most} of what we can measure from a transformer.
Yet leveraging these measurements to predict and control behavior remains difficult.
For example, in an effort to both detect and mitigate social biases (e.g., gender bias)
% \footnote{For example, they detect gender bias by observing how Claude fills
% pronouns in ambiguous contexts: "The boss explained a new accounting system to the 
% administrative assistant, who then taught the rest of \_ team."}
in Claude 3 Sonnet,
\citet{durmus2024steering} analyzed its representations.
Their results were mixed. It is possible to detect
certain social biases from representations. But intervening on these
representations can impair general behavior---a result consistent with
the systematic study in \citet{da2025steering}. In short, our tools for understanding representations
are not reliable enough to predict or control behavior.
%%%%%% Subfield gap %%%%%
% There is a history of validating methods by verifying against
% ground-truth.
% - Could a neuroscientist understand a microprocessor?
% Have some tools for doing this in understanding reps (TRACR, Alta)

To address this, \citet{lindner2023tracr} proposed compilers as a tool for studying representations. 
By compiling known algorithms into a transformer, techniques for understanding their representations
can be verified against a ground truth. Consider the conditional string reversal algorithm 
in Fig. \ref{fig:conditional_reverse}.

\begin{wrapfigure}{l}{0.62\linewidth}
\vspace{-.8em}
\includegraphics[width=\linewidth]{figures/conditional_reverse/conditional_reverse.pdf}
\captionof{figure}{Reversing a string if its first character is a vowel.}
\label{fig:conditional_reverse}
\vspace{-1.0em}
\end{wrapfigure}

If we could compile this algorithm\footnote{The code is Python with parallel dictionary access. Passing a list
of keys to a dictionary returns a list of values.} into a transformer, 
we would have a rich model of the underlying representations, and can use that ground truth 
to verify whether techniques for understanding representations can recover it.
Similar metascientific studies have been explored in other sciences: \citet{jonas2017could} 
asked whether neuroscientists could understand a microprocessor, and \citet{lazebnik2002can}
whether a biologist could fix a radio.

\begin{wrapfigure}{r}{0.40\linewidth}
  \includegraphics[width=\linewidth]{figures/transformer/transformer.pdf}
  \vspace{.1em}
  \captionof{figure}{Compiling program $\cm{e}$ into an attention head of a transformer.}
  \label{fig:transformer}
  \vspace{-4em}
\end{wrapfigure}

%%%%% Our gap %%%%%
% But the closed program assumption sucks!!!!!!!!!!
However, compiling algorithms into transformers has so far been of limited use in verifying our 
understanding of representations.
In part, this is because prior compilers make a \textit{closed program} assumption 
\citep{perconti2014verifying}. They cannot compile programs with free variables. As a 
consequence, the program must specify the entire algorithm a transformer implements. 
This restricts their application to hand-constructed transformers rather than ones trained 
by gradient descent---which are ultimately the kind we want to understand.

%%%%% Contributions %%%%%
% Here is how we fix it.
\paragraph{Contributions} This work resolves that limitation with $\cj$\!, a language whose
programs compile to transformers which can be subsequently trained by gradient descent. For example,
the conditional string reversal program in Fig. \ref{fig:conditional_reverse} can be encoded in
$\cj$\! and then compiled into an attention head of a transformer trained to solve the task,
as shown in Fig. \ref{fig:transformer}. More generally, our contributions are:

\setlist[itemize]{topsep=4pt, itemsep=2pt, parsep=2pt, leftmargin=1.5em, rightmargin=1.5em}
\begin{itemize}[label=$\bullet$]
\item We specify the syntax, evaluator, and compiler for $\cj$\! programs (Section \ref{sec:cajal}).
\item We prove that $\cj$\! programs compile correctly to transformers (Section \ref{sec:cajal}).
\item We implement $\cj$\!, programming
a transformer to solve a conditional string reversal task using a compiled attention head (Section \ref{sec:experiments})
\item We analyze its representations using principal components analysis (PCA),
 ablations, and steering to confirm it uses the compiled attention head
 to solve the task (Section \ref{sec:experiments}).
\end{itemize}

As a result, $\cj$\! enables verification of interpretability techniques against partial ground truths
in transformers trained by gradient descent.

% TODO:
% 1. Fix section references now. Put stubs for section references, also for the appendices. This 
% way we don't have to go back and fix them later!




